\section{Discussion and Conclusion}

\noindent\textbf{Main findings.} MeshTac turns unlabeled real contacts into geometric supervision by aligning a known object mesh to each image. The ablations connect the improvement to the quality of this alignment: refining the contact state improves the projected labels, and the student trained on those labels reconstructs contact more accurately. Because the gain remains when synthetic images are removed, the evidence points to contact-state error---rather than appearance synthesis alone---as the main obstacle to using unlabeled contacts within a capture session.

\noindent\textbf{Lessons for label-efficient tactile learning.} Geometry coverage and sensor appearance play different roles. Real unlabeled images already provide the target appearance, so mesh alignment is a strong first step when registered contacts are scarce. Synthetic data becomes more useful when coverage changes, as seen in the second session and the unseen-object pose study. This suggests a practical order of operations: first establish a renderer-free alignment baseline, then add simulation when pose coverage or session robustness is limiting. Comparisons should disclose the real data used by the renderer, teacher and aligner, together with validation labels used for checkpoint selection; otherwise renderer supervision can be mistaken for a gain from simulation. Image fidelity can screen renderers, but downstream contact metrics remain necessary because lower all-pixel error need not mean better contact localization.

\noindent\textbf{Limitations and future work.} MeshTac assumes known objects: the predictor uses object identity, while pseudo-labeling and optional refinement require a mesh. The targets are mesh projections, so the experiments measure registered object geometry rather than physical elastomer deformation. The 12 ridge-pattern objects, labeled validation set and quasi-static presses leave curved and textured contacts, unknown objects, shear and slip untested. Repeated ridges also create translation ambiguity, and a small set of cross-session alignments retain large errors. Future work should develop uncertainty-aware or symmetry-aware alignment, extend the labels to dynamic contact, and test whether the geometric gains improve closed-loop manipulation. Relaxing the known-mesh assumption through joint shape and contact-state estimation is a further direction.
